\section{Introducción: no es un perfil biográfico, sino una línea de investigación}

Esta entrevista dura seis horas y cuarenta y cinco minutos. En apariencia, es la primera vez que Xie Saining (谢赛宁) relata de forma sistemática su formación, sus estudios, su investigación, su emprendimiento y sus criterios de valor. El hilo más profundo es otro: cómo un investigador que trabaja en computer vision y representation learning redefine «visión», «modelo del mundo» y «cerebro robótico» en un momento en que la narrativa de los LLM se impone sobre todo lo demás. Por eso, esta nota no repetirá la entrevista palabra por palabra como un relato cronológico, sino que la organiza como una línea de investigación que se puede estudiar.

Este episodio tiene tres ejes. El primero es la trayectoria personal: de la clase ACM de la Universidad Jiao Tong de Shanghái a UCSD, FAIR, DeepMind, NYU y AMI Labs; cada decisión giró en torno a «con quién trabajar, qué problema abordar y si hay recursos suficientes». El segundo es la línea técnica: de las tareas de visión, las representaciones jerárquicas y la comprensión de video al predictive world model. El tercero es la línea organizacional: cuando un problema de investigación requiere recursos muy superiores a los de un laboratorio orientado a publicar artículos, pero no conviene que la carrera armamentista de productos lo absorba por completo, ¿cómo puede una empresa emergente conservar el research oxygen?

\begin{importantbox}{Tesis central del episodio}
Xie Saining insiste a lo largo de la entrevista en que los LLM son un componente importante de los sistemas inteligentes, pero no son todo. Una inteligencia orientada de verdad al mundo físico necesita aprender representaciones relevantes para la tarea a partir de señales continuas, de alta dimensión y con ruido, predecir las consecuencias de una action y usar esas predicciones para planning, decision making y control de seguridad.
\end{importantbox}

\begin{knowledgebox}{Nota sobre la estrategia visual}
Este video muestra una toma fija de entrevista, sin slides didácticas, pizarrón, demostraciones de producto ni gráficas legibles. Conforme al estándar de pódcast de este repositorio, el cuerpo del texto no repite fotogramas de los participantes; la portada sirve para identificar la fuente, y el contenido didáctico se presenta mediante el ciclo del modelo del mundo, la escalera de capacidades, el modelo organizacional y tablas de conceptos.
\end{knowledgebox}

\subsection{Resumen de la sección}

El valor de este episodio no está en saber a quién rechazó Xie Saining o dónde trabajó, sino en entender cómo una postura técnica se forma poco a poco a partir de la experiencia personal, el criterio de investigación y las condiciones organizacionales. Las secciones siguientes dividen la entrevista en: la trayectoria de investigación, la definición de modelo del mundo, el límite entre los LLM y el mundo físico, la propuesta organizacional de AMI y, al final, la reflexión sobre la inteligencia y las citas filosóficas.

